\documentclass[12pt]{article}
\usepackage[T1]{fontenc}
\usepackage[utf8]{inputenc}
\usepackage{lmodern}
\usepackage{textcomp}
\usepackage{enumitem}
\usepackage{geometry}
\geometry{margin=1in}
\begin{document}
\begin{center}
Answer Key - Economics - Undergraduate Level Assessment 8 \
\small (Answers are ordered exactly as in the question paper.)
\end{center}

\section*{Multiple Choice}
\begin{enumerate}[label=\arabic*.]
\item \textbf{Answer:} E
\\ \textit{Options:} A) Unions determine labor productivity; B) Unions can directly control employment rates nationally; C) Unions are the primary drivers of economic growth and thereby increase labor demand; D) Unions can significantly increase labor demand; E) Unions have limited means to influence the demand for labor
\\ \textit{Note:} Unions primarily influence wages and working conditions rather than directly affecting the overall demand for labor, which is largely determined by market forces and the broader economy.
\item \textbf{Answer:} A
\\ \textit{Options:} A) decreased government spending and increased taxes.; B) the federal government to decrease the money supply.; C) the federal government to run a deficit.; D) decreased taxes and increased government spending.; E) the federal government to increase the money supply.
\\ \textit{Note:} Decreased government spending and increased taxes are appropriate fiscal policies to remedy inflation because they reduce overall demand in the economy, helping to lower price levels by curbing excessive consumer spending and investment.
\item \textbf{Answer:} A
\\ \textit{Options:} A) If both the demand and supply curves are price inelastic; B) If the demand curve is perfectly elastic and the supply curve is perfectly inelastic; C) If the demand curve is price inelastic and the supply curve is price elastic; D) If both the demand and supply curves are price elastic; E) If the demand curve is price elastic and the supply curve is price inelastic
\\ \textit{Note:} Consumers would bear the largest burden of an excise tax when both the demand and supply curves are price inelastic because this means that neither consumers nor producers can easily change their quantity demanded or supplied in response to price changes, leading to a greater transfer of the tax burden to consumers.
\item \textbf{Answer:} E
\\ \textit{Options:} A) the cost of all produced units; B) the fixed cost associated with the production facility; C) the cost of total production; D) the cost of initial unit of production; E) the additional cost of one more unit of production
\\ \textit{Note:} The term marginal costs refers to the incremental expense incurred when producing an additional unit of a good or service, highlighting the relationship between production levels and cost efficiency in economics.
\item \textbf{Answer:} D
\\ \textit{Options:} A) Very few competitors.; B) One firm sets the market price.; C) Identical products offered by each firm.; D) Excess capacity.; E) Significant barriers to entry.
\\ \textit{Note:} Excess capacity in monopolistic competition occurs because firms produce below their optimal output level due to product differentiation and the presence of many competitors, leading to inefficiencies in resource allocation.
\end{enumerate}

\section*{True / False}
\begin{enumerate}[label=\arabic*.]
\setcounter{enumi}{5}
\item \textbf{Answer:} True
\\ \textit{Options:} A) True; B) False
\\ \textit{Note:} The statement is true because a fitted regression line can include polynomial terms, such as \(x_t^2\), allowing for non-linear relationships between the independent variable and the dependent variable, which is captured by the estimated coefficients \(̂α\) and \(̂β\) along with the error term \(̂u_t\).
\item \textbf{Answer:} True
\\ \textit{Options:} A) True; B) False
\\ \textit{Note:} A monopoly faces a downward-sloping demand curve, which means that to sell additional units, it must lower the price, leading to a marginal revenue curve that is initially downward sloping as the revenue gained from selling one more unit decreases with each additional unit sold.
\item \textbf{Answer:} False
\\ \textit{Options:} A) True; B) False
\\ \textit{Note:} The fixed effects panel model and the random effects model are distinct econometric approaches that handle unobserved heterogeneity differently, with fixed effects controlling for time-invariant characteristics while random effects assumes these characteristics are uncorrelated with the explanatory variables.
\item \textbf{Answer:} False
\\ \textit{Options:} A) True; B) False
\\ \textit{Note:} A decrease in the cost of labor for bookbinding reduces production costs, which would actually shift the supply curve for textbooks to the right, not to the left.
\item \textbf{Answer:} False
\\ \textit{Options:} A) True; B) False
\\ \textit{Note:} A nation producing inside its production possibility curve indicates inefficiency or underutilization of resources, but it does not imply a permanent decrease in production capacity, as the nation could improve efficiency or reallocate resources to reach its potential output.
\end{enumerate}

\section*{Long Form Response}
\begin{enumerate}[label=\arabic*.]
\setcounter{enumi}{10}
\item \textbf{Answer:} A wage rate of \$4.50 in a market economy has significant implications for labor hiring decisions, particularly concerning the supply and demand for labor. Employers will assess whether the productivity of workers at this wage meets or exceeds the cost of hiring them. To determine the number of labor units hired at this wage, one would analyze the labor demand curve, which reflects the number of workers firms are willing to employ at various wage levels. If the wage is below the market equilibrium, firms may hire more workers due to lower labor costs, while higher wages could lead to a contraction in hiring as firms seek to minimize expenses. Ultimately, the intersection of the labor demand curve with the wage rate will indicate the quantity of labor units hired.
\\ \textit{Note:} correct because it identifies the relationship between wage rates, labor supply and demand, and emphasizes the importance of analyzing the labor demand curve to determine the number of workers hired at a given wage, particularly when the wage is below the market equilibrium.
\item \textbf{Answer:} Pure time-series models, while useful for forecasting and identifying patterns in economic data, have significant disadvantages due to their lack of theoretical motivation. Unlike structural models, which are grounded in economic theory and incorporate underlying relationships between variables, time-series models often rely solely on historical data without explaining the mechanisms behind observed trends. This reliance can lead to misleading forecasts if the underlying economic conditions change, as the model does not account for shifts in structural relationships. Additionally, the absence of theoretical foundations can limit the model's applicability in policy analysis, making it difficult to draw meaningful conclusions about causal relationships or the impact of economic interventions. Ultimately, the lack of theoretical motivation in pure time-series models undermines their effectiveness in providing comprehensive insights into economic phenomena.
\\ \textit{Note:} The answer correctly highlights that pure time-series models lack theoretical foundations, making them less reliable for understanding economic dynamics, as they focus on historical data patterns without addressing the causal relationships that structural models provide.
\end{enumerate}

\vfill
\noindent\textit{End of Answer Key}
\end{document}